\section{Geometric and geodesic consequences}
\label{sec:consequences}

The first two subsections concern the exponential law of arbitrary rate
$\lambda>0$. Theorem~\ref{thm:exponential-shape} supplies both regularity
properties needed there. The final two subsections apply
Theorem~\ref{thm:gamma-differentiability} to a quantitative midpoint
estimate and to cone-tip competition for the whole Gamma family.

For a Euclidean unit vector $u$, put $v_u=u/\mu(u)$ and
$\rho_u=\nabla\mu(u)$. Homogeneity gives
$\rho_u\cdot v_u=1$, and the supporting inequality gives
$\rho_u\cdot x\le\mu(x)$ for every $x$. Hence
\[
 P_u=\{x:\rho_u\cdot x=1\}
\]
is the supporting line at $v_u$. Strict convexity gives
$P_u\cap\Bcal=\{v_u\}$: two contact points would put their entire
segment in the boundary. Thus the set of radial directions of contact
with $P_u$ consists of the single direction $u$.

\subsection{A fixed deterministic direction}

A finite geodesic is a path realizing the passage time between its
endpoints. An infinite geodesic, or a bigeodesic, is respectively a
one-sided or two-sided infinite path every finite segment of which is
a finite geodesic. For the Gamma laws in this paper, finite geodesics
exist and are almost surely unique. Indeed, choose $\theta>0$ so large
that $3\E e^{-\theta\tau_e}<1$. Counting simple paths shows that the
probability of a simple path from a fixed vertex with at least $m$
edges and cost at most $t$ is bounded by
$C e^{\theta t}(3\E e^{-\theta\tau_e})^m$. Thus almost surely only
finitely many simple paths from that vertex have cost below any fixed
bound. Thus every passage ball $\{z:T(x,z)\le t\}$ is finite, and a
minimum is attained between any two vertices. Two distinct simple paths with the
same endpoints have different edge sets; conditioning on all but an
edge in their symmetric difference shows that their costs are unequal
almost surely. Countability gives uniqueness simultaneously for all
lattice endpoints.

A sequence $(z_n)$ is directed toward a unit vector $u$ when
$\lVert z_n\rVert_2\to\infty$ and
$z_n/\lVert z_n\rVert_2\to u$. An infinite geodesic has direction $u$
when its vertex sequence has this property. Write $\Gamma(x,y)$ for the
unique finite geodesic between lattice vertices $x$ and $y$. Local
convergence of paths rooted at $x$ means eventual agreement of each fixed
finite initial segment.

\begin{corollary}\label{cor:fixed-direction}
Fix a deterministic unit vector $u$. There is an event $\Omega_u$ of
probability one on which all the following statements hold.
\begin{enumerate}
\item For every $x\in\Z^2$ there is exactly one infinite geodesic
      $\Gamma_x^u$ from $x$ with direction $u$. For every sequence $(z_n)$
      directed toward $u$, including a sequence chosen from the edge
      configuration,
      $\Gamma(x,z_n)\to\Gamma_x^u$ locally.
\item For every $x,y\in\Z^2$, the paths $\Gamma_x^u$ and
      $\Gamma_y^u$ have a common tail.
\item No bigeodesic has an end with direction $u$.
\item For every $x,y\in\Z^2$, the limit
\[
 B_u(x,y)=\lim_{n\to\infty}\bigl(T(x,z_n)-T(y,z_n)\bigr)
\]
      exists and is independent of the choice of sequence directed toward
      $u$. Moreover,
\[
 \E B_u(0,x)=\nabla\mu(u)\cdot x,
\]
      and, for every $\varepsilon>0$, only finitely many $x\in\Z^2$
      satisfy
\[
 |B_u(0,x)-\nabla\mu(u)\cdot x|>\varepsilon\lVert x\rVert_1.
\]
\end{enumerate}
\end{corollary}

\begin{proof}
Damron--Hanson's product-measure hypothesis A1$'$ holds: the common law
is continuous, and the minimum $Y$ of the four weights incident to the
origin is exponential of rate $4\lambda$, whence
$\E Y^2=1/(8\lambda^2)<\infty$.
Their geometric assumptions (2.1)--(2.2) require differentiability at
the chosen boundary point and at both endpoints of its supporting face.
The preceding supporting-line calculation identifies that face with
$\{v_u\}$. Thus their supporting sector is exactly the direction $u$.
Theorems~1 and~2 of \cite{DH2017} give the convergence, coalescence,
bigeodesic exclusion, and Busemann statements on one probability-one
event for this fixed $u$. Their normalized supporting vector is
$\rho_u=\nabla\mu(u)$ by that calculation.

For completeness, uniqueness follows from the convergence assertion.
If $(x=x_0,x_1,\ldots)$ is any infinite geodesic of direction $u$, then
its segment ending at $x_n$ is $\Gamma(x,x_n)$ by uniqueness of finite
geodesics. These segments converge both to that infinite geodesic and
to $\Gamma_x^u$, so the paths agree.
\end{proof}

The Busemann limit also has an elementary pathwise description. If $c$
is any common vertex after $\Gamma_x^u$ and $\Gamma_y^u$ have coalesced,
local convergence implies that both finite geodesics to $z_n$ pass
through $c$ for all sufficiently large $n$. Hence
\begin{equation}\label{eq:busemann-common-vertex}
 B_u(x,y)=T(x,c)-T(y,c).
\end{equation}
In particular the displayed difference defining $B_u(x,y)$ is eventually
constant for each fixed pair and each such target sequence. Taking limits
in the triangle inequality and in differences of passage times gives
\[
 |B_u(x,y)|\leq T(x,y),\qquad
 B_u(x,z)=B_u(x,y)+B_u(y,z).
\]
The first inequality also gives integrability. If $y$ lies on
$\Gamma_x^u$, then $B_u(x,y)=T(x,y)$. Translation of the environment and
of the targets shows that this additive random function is stationary.

The direction $u$ is fixed before taking the probability-one event.
One may intersect these events for a prescribed countable set of
directions. The result does not assert simultaneous uniqueness in all
directions, or simultaneous absence of all bigeodesics.

Independently of these shape arguments, the companion manuscript
\cite[Theorem~1.1]{OpenAINoBigeodesics2026} excludes all bigeodesics on one
probability-one event for iid $\operatorname{Gamma}(\kappa,\lambda)$ edge
weights, for each fixed $\kappa,\lambda>0$. This includes the exponential
laws here; these distributions satisfy its nonatomic minimum-second-moment
hypotheses.

\subsection{Passage through a prescribed vertex}

\begin{corollary}\label{cor:midpoint}
For every nonzero $v\in\Z^2$,
\[
 \lim_{n\to\infty}
 \PP\bigl(\lfloor n/2\rfloor v\in\Gamma(0,nv)\bigr)=0.
\]
More generally, let $a_n,b_n,z_n$ be deterministic lattice vertices with
\[
 \lVert a_n-z_n\rVert_2\to\infty,
 \quad\lVert b_n-z_n\rVert_2\to\infty,
 \quad
 \frac{b_n-z_n}{\lVert b_n-z_n\rVert_2}\longrightarrow u
\]
for a deterministic unit vector $u$. Then
$\PP(z_n\in\Gamma(a_n,b_n))\to0$.
\end{corollary}

\begin{proof}
Set $p_n=a_n-z_n$ and $q_n=b_n-z_n$. Translation invariance gives
$\PP(z_n\in\Gamma(a_n,b_n))=
\PP(0\in\Gamma(p_n,q_n))$.
If the asserted probabilities had positive limsup,
the event that $0\in\Gamma(p_n,q_n)$ for infinitely many $n$ would
have positive probability. Intersect this event with $\Omega_u$ from
Corollary~\ref{cor:fixed-direction}.
Restrict to the infinite set of indices for which
$0\in\Gamma(p_n,q_n)$. On the intersection, the arm from $0$ to $q_n$ converges locally to
$\Gamma_0^u$. Both arms have lengths tending to infinity, since their
endpoints tend to infinity. Local finiteness of the lattice permits a
diagonal subsequence along which the other arm also converges to an
infinite path. Every finite segment of the resulting two-sided path is
a segment of one of the finite geodesics $\Gamma(p_n,q_n)$, and hence
is geodesic. The two arms have no vertex in common except $0$, because
each finite geodesic is simple. This constructs a bigeodesic with an
end of direction $u$, contradicting Corollary~\ref{cor:fixed-direction}.
The first assertion uses $a_n=0$, $b_n=nv$, and
$z_n=\lfloor n/2\rfloor v$.
\end{proof}

This is the compactness argument underlying \cite[Theorem~2.1]{DH2017}.
Ahlberg and Hoffman \cite[Theorem~2.1]{AH2019} already prove the whole
qualitative corollary for arbitrary deterministic endpoint sequences
escaping from the tested site, assuming continuity and the
minimum-of-four second moment. This includes all Gamma laws here and
requires no shape regularity. The proof above records its connection
with the fixed-direction consequences.

\subsection{A quantitative midpoint bound for Gamma weights}

Differentiability has a useful consequence even when strict convexity
has not been established: a $C^1$ planar convex body cannot have only
finitely many extreme points. This verifies the geometric hypothesis
of the quantitative midpoint theorem of Dembin, Elboim, and Peled.

\begin{corollary}[Quantitative midpoint bound for the Gamma family]
\label{cor:gamma-quantitative-midpoint}
For every $\kappa>0$ there is $C_\kappa<\infty$ such that, for every
$\lambda>0$ and all $a,b,z\in\Z^2$ with
$D=\min\{\lVert a-z\rVert_1,\lVert b-z\rVert_1\}\geq1$,
\[
 \PP\bigl(z\in\Gamma_{\kappa,\lambda}(a,b)\bigr)
 \leq C_\kappa\frac{\log^2(D+2)}{D^{1/16}},
\]
where $\Gamma_{\kappa,\lambda}$ denotes the unique geodesic for
iid $\operatorname{Gamma}(\kappa,\lambda)$ edge weights.
\end{corollary}

\begin{proof}
Theorem~\ref{thm:gamma-differentiability} gives a $C^1$ boundary. We verify that such a
compact planar convex body has infinitely many extreme points.
At a boundary point choose a supporting line. Its intersection with
the body is a point or a compact segment, and its endpoints are
extreme points of the body: any convex decomposition of an endpoint
must stay in the supporting line, where extremality is immediate.
Thus every boundary point is a convex combination of extreme points.
Every interior point lies between two boundary points, so the whole
body is the convex hull of its extreme points. If there were only
finitely many, the body would be a polygon with nonempty interior.
At a vertex of that polygon the two adjacent sides give distinct
supporting lines, contradicting differentiability. Hence the number
of extreme points exceeds forty.

The Gamma density is absolutely continuous, has no atom at zero,
and has a finite exponential moment by Lemma~\ref{df:ga:moment}.
All the hypotheses of \cite[Theorem~1.2]{DEP2024} are therefore
satisfied. Apply that theorem at rate one to obtain $C_\kappa$.
The coupling $t_e^{(\kappa,\lambda)}=
t_e^{(\kappa,1)}/\lambda$ leaves every geodesic unchanged and proves
the assertion for all positive rates with the same constant.
\end{proof}

For Gamma laws, differentiability also verifies the tangent and
supporting-face endpoint assumptions of \cite{DH2017}. Its conclusions
then retain the supporting-sector formulation. The single-direction
statements of Corollary~\ref{cor:fixed-direction} use the strict
convexity established for exponential weights.

\subsection{Cone-tip competition for Gamma weights}

We apply the tangent theorem of Ahlberg, Deijfen, and Sfragara
\cite[Theorem~1]{ADS2026} to their simultaneous cone-versus-tip initial
configuration.

\begin{corollary}[Gamma and Richardson cone-tip survival]
\label{cor:gamma-cone-tip}
Fix $\kappa,\lambda>0$, a deterministic unit vector $\theta\in S^1$,
and a geometric half-angle $0\le\alpha\le\pi$. Define the closed cone
\[
 \begin{aligned}
 C^\theta_\alpha&=\{0\}\cup
 \{x\in\R^2\setminus\{0\}:\angle(x,-\theta)\le\alpha\},\\
 I_1&=(C^\theta_\alpha\cap\Z^2)\setminus\{0\},
 \qquad I_2=\{0\},
 \end{aligned}
\]
where $\angle(x,-\theta)$ is the Euclidean angle between the two vectors.
Use one iid $\operatorname{Gamma}(\kappa,\lambda)$ field on unoriented
nearest-neighbor edges of $\Z^2$, with passage metric $T$. At time zero,
start type 1 from every site of $I_1$ and type 2 from $I_2$.
Both types use this same realization of $T$, and each site is permanently
claimed by the type that reaches it first. With
$T(I,z)=\inf_{x\in I}T(x,z)$, interpreting the infimum of the empty set
as $\infty$, the sites eventually claimed by type 2 are almost surely
\[
 \mathcal E_2=\{z\in\Z^2:T(0,z)<T(I_1,z)\}.
\]
Then
\[
 \PP\bigl(|\mathcal E_2|=\infty\bigr)>0
 \quad\Longleftrightarrow\quad
 \alpha<\pi/2.
\]
In particular, at $\alpha=\pi/2$, when the cone is the closed half-plane
$\{x:x\cdot\theta\le0\}$, type 2 claims only finitely many sites almost
surely. At $\kappa=1$, this gives Benjamini's cone-tip prediction for
equal-rate Richardson competition with any common rate $\lambda>0$ in
the same initial configuration.
\end{corollary}

\begin{proof}
The Gamma law is atomless for every $\kappa>0$, including $\kappa<1$.
For four independent samples,
\[
 \E\bigl[\min(\tau_1,\ldots,\tau_4)^2\bigr]
 \le \E\tau_1^2
 =\frac{\kappa(\kappa+1)}{\lambda^2}<\infty.
\]
Thus the continuity and minimum-second-moment hypotheses of
\cite[Theorem~1]{ADS2026} hold.

The infinite seed set causes no ambiguity in the arrival-time comparison.
If $I_1$ is nonempty, choose one seed as a finite upper bound for
$T(I_1,z)$. Only finitely many seeds lie in the corresponding passage
ball about $z$, so the infimum is attained. Simple paths from distinct
seeds to $z$ have different edge sets. Conditioning on all edge weights
except one in their symmetric difference and using atomlessness rules
out equal costs;
countability gives this simultaneously for all such paths and sites.
Thus the shared-metric competition is described by the displayed strict
comparison defining $\mathcal E_2$.

Theorem~\ref{thm:gamma-differentiability} gives a $C^1$ boundary.
For the fixed $\theta$, put
$h(\theta)=\max_{x\in\Bcal}x\cdot\theta>0$. The supporting line
\[
 L_\theta=\{x:x\cdot\theta=h(\theta)\}
\]
is perpendicular to the positive $\theta$-ray and meets that ray at
$h(\theta)\theta$. At any contact point with $\Bcal$, the $C^1$ boundary
makes $L_\theta$ the unique supporting line, hence a tangent in the
sense of \cite{ADS2026}; the contact point need not lie on the ray.
For a seed set $I\subset\Z^2\setminus\{0\}$, write
$\mathcal E_2(I)=\{z\in\Z^2:T(0,z)<T(I,z)\}$, with the same
empty-set convention for $T(I,z)$. If $I\subset J$, then
$T(J,z)\le T(I,z)$ and hence
$\mathcal E_2(J)\subset\mathcal E_2(I)$. If $\alpha<\pi/2$, choose
$\alpha<\alpha'<\pi/2$. The source's cone of half-angle $\alpha'$
contains the closed $\alpha$-cone, regardless of its boundary convention,
so its Theorem~1 gives positive survival for the closed $\alpha$-cone
as well. If
$\alpha\ge\pi/2$, the closed $\alpha$-cone contains the closed critical
half-plane. Theorem~5 gives almost-sure finite type-2 growth for that
half-plane, and monotonicity gives the same conclusion for the larger
cone. This proves the stated closed-cone convention throughout
$[0,\pi]$, without strict convexity, for every prescribed orientation.

When $\kappa=1$, the shared exponential edge field gives the equal-rate
two-type Richardson model. The coupling
$\tau_e^{(\lambda)}=\tau_e^{(1)}/\lambda$ multiplies both arrival
times by the same positive constant and leaves every comparison
defining $\mathcal E_2$ unchanged. This proves the stated specialization.
\end{proof}

The probability assertion is for each fixed orientation, not a simultaneous
pathwise statement over all directions. Below the threshold it asserts a
positive survival probability without claiming that this probability is one.
It does not address unequal rates, independent type-specific metrics,
coexistence from two finite seed sets, or the rate or asymptotic shape of
type-2 growth.
